\documentclass[11pt]{article}
\usepackage[a4paper,margin=1in]{geometry}
\usepackage{amsmath,amssymb,amsthm,mathtools}
\usepackage{stmaryrd}
\usepackage{mathpartir}
\usepackage{hyperref}
\usepackage[T1]{fontenc}
\usepackage[scaled=0.92]{inconsolata}

\newtheorem{theorem}{Theorem}
\newtheorem{lemma}[theorem]{Lemma}
\newtheorem{proposition}[theorem]{Proposition}
\newtheorem{corollary}[theorem]{Corollary}
\newtheorem{definition}[theorem]{Definition}

% --- Notation ---
\newcommand{\D}{\mathsf{D}}
\newcommand{\dirac}{\mathsf{dirac}}
\newcommand{\bind}{\mathbin{\gg\!=}}
\newcommand{\sem}[2][]{\llbracket #2 \rrbracket_{#1}}
\newcommand{\E}[1]{E_{#1}}
\newcommand{\ty}[1]{\mathsf{#1}}
\newcommand{\bool}{\ty{bool}}
\newcommand{\num}{\ty{num}}
\newcommand{\arr}{\rightarrow}
\newcommand{\val}{\mathsf{val}}
\newcommand{\ul}[1]{\underline{#1}} % numeral for a real
\newcommand{\FV}{\mathsf{FV}}
\newcommand{\holes}{\mathsf{holes}}
\newcommand{\N}{\mathbb{N}}
\newcommand{\R}{\mathbb{R}}
\newcommand{\RR}{\mathbb{R}}
\newcommand{\B}{\mathcal{B}}
\newcommand{\Tt}{\mathcal{T}}
\newcommand{\plug}{\mathsf{plug}}
\newcommand{\shape}{\mathsf{shape}}

\title{A Small-Step Semantics for a Tiny Probabilistic Language}
\author{}
\date{}

\begin{document}
\maketitle

\paragraph{Syntax.}
\[
\begin{array}{rrcl}
\text{Types} & t &::=& \num \mid \bool \mid t_1 \arr t_2
\\[1ex]
\text{Terms} & e &::=& x \mid \lambda x{:}t.\,e \mid e_1\,e_2 \mid \ty{if}\;e_1\;\ty{then}\;e_2\;\ty{else}\;e_3 
\\
&&\mid& e_1{<}e_2 \mid \ul{r} \ \ (r\in\RR) \mid \ty{true} \mid \ty{false} \mid \ty{unif}(a,b)\ \ (a<b\in\RR)
\\[1ex]
\text{Values} & v &::=& \ty{true}\mid \ty{false}\mid \lambda x{:}t.\,e \mid \ul{r} \ \ (r\in\RR)
\end{array}
\]

\paragraph{Typing.} Typing judgments have the form $\Gamma\vdash e:t$.
We write $\,\E{t} \triangleq \{\,e\mid \emptyset\vdash e:t\,\}\,$ for closed well-typed terms.

\begin{mathpar}
\inferrule*[Right=T-Var]{x{:}t\in\Gamma}{\Gamma\vdash x:t}
\and
\inferrule*[Right=T-Num]{\ }{\,\Gamma\vdash \ul{r}:\num\,}
\and
\inferrule*[Right=T-BoolT]{\ }{\,\Gamma\vdash \ty{true}:\bool\,}
\\
\inferrule*[Right=T-BoolF]{\ }{\,\Gamma\vdash \ty{false}:\bool\,}
\and
\inferrule*[Right=T-Unif]{a<b}{\,\Gamma\vdash \ty{unif}(a,b):\num\,}
\\
\inferrule*[Right=T-Lam]{\Gamma,x{:}t\vdash e:u}{\Gamma\vdash \lambda x{:}t.\,e : t\arr u}
\and
\inferrule*[Right=T-App]{\Gamma\vdash e_1:t\arr u \\ \Gamma\vdash e_2:t}{\Gamma\vdash e_1\,e_2:u}
\\
\inferrule*[Right=T-If]{\Gamma\vdash e_1:\bool \\ \Gamma\vdash e_2:t \\ \Gamma\vdash e_3:t}{\Gamma\vdash \ty{if}\;e_1\;\ty{then}\;e_2\;\ty{else}\;e_3:t}
\and
\inferrule*[Right=T-Lt]{\Gamma\vdash e_1:\num \\ \Gamma\vdash e_2:\num}{\Gamma\vdash e_1<e_2:\bool}
\end{mathpar}

\paragraph{Values.} The predicate $\val(e)$ holds exactly for $v$ in the grammar above.

\section*{Small-step semantics}



For each $t$ we define a function $\sem[t]{\cdot}:\E{t}\to \D(\E{t})$ by structural equations,
where $\D(\E{t})$ is the space of probability measures on $\E{t}$.
We explain the measurable structure on expressions in the next section.
We write $\sem{e}$ when $t$ is clear.

\paragraph{Distribution interface.}
We work abstractly with a distribution functor $\D(-)$ equipped with
$\dirac:A\to \D(A)$ and bind $\bind:\D(A)\to (A\to \D(B))\to \D(B)$.
We write $U[a,b]\in \D([a,b])$ for the uniform distribution on the interval.

\paragraph{Values and randomness.}
\begin{align*}
\sem[\num]{\ul r} &= \dirac(\ul r)
&
\sem[\bool]{\ty{true}} &= \dirac(\ty{true})
&
\sem[\bool]{\ty{false}} &= \dirac(\ty{false})
\\
\sem[{t\arr u}]{\lambda x{:}t.\,e} &= \dirac(\lambda x{:}t.\,e)
&
\sem[\num]{\ty{unif}(a,b)} &= \iota_{*}(U[a,b])
\end{align*}
where $\iota:[a,b]\to \E{\num}$ is $r\mapsto \ul r$ and $\iota_{*}$ denotes pushforward.

\paragraph{Conditionals.}
\[
\sem{\,\ty{if}\;e_1\;\ty{then}\;e_2\;\ty{else}\;e_3\,}
=
\begin{cases}
\dirac(e_2) & \text{if } e_1=\ty{true}\\
\dirac(e_3) & \text{if } e_1=\ty{false}\\
\sem{e_1}\bind (\lambda e_1'.\, \dirac(\ty{if}\;e_1'\;\ty{then}\;e_2\;\ty{else}\;e_3)) & \text{if } \neg\val(e_1)
\end{cases}
\]

\paragraph{Comparison.}
\[
\sem{\,e_1{<}e_2\,}
=
\begin{cases}
\dirac(\ty{true})  & \text{if } e_1=\ul r_1,\ e_2=\ul r_2,\ r_1<r_2\\
\dirac(\ty{false}) & \text{if } e_1=\ul r_1,\ e_2=\ul r_2,\ r_2\le r_1\\[0.4ex]
\sem{e_1}\bind (\lambda e_1'.\, \dirac(e_1'{<}e_2)) & \text{if } \neg\val(e_1)\\
\sem{e_2}\bind (\lambda e_2'.\, \dirac(e_1{<}e_2')) & \text{if } \val(e_1)\wedge \neg\val(e_2)
\end{cases}
\]

\paragraph{Application (left-to-right CBV).}
\[
\sem{\,e_1\,e_2\,}
=
\begin{cases}
\dirac(e[v/x]) & \text{if } e_1=\lambda x{:}t.\,e \text{ and } \val(e_2)=v\\
\sem{e_1}\bind (\lambda e_1'.\, \dirac(e_1'\,e_2)) & \text{if } \neg\val(e_1)\\
\sem{e_2}\bind (\lambda e_2'.\, \dirac(e_1\,e_2')) & \text{if } \val(e_1)\wedge \neg\val(e_2)
\end{cases}
\]

\bigskip
\noindent
(Variables never appear in $\E{t}$, so no clause for $x$.)
\\[1em]
We shall now prove that this defines a function $\sem[t]{\cdot}:\E{t}\to \D(\E{t})$.

\section*{Typing lemmas}

\begin{lemma}[Substitution]\label{lem:subst}
If $\Gamma,x{:}t\vdash e:u$ and $\emptyset\vdash v:t$ then $\Gamma\vdash e[v/x]:u$.
\end{lemma}
\begin{proof}
By a routine induction on the derivation of $\Gamma,x{:}t\vdash e:u$ and a weakening lemma.
\end{proof}

\begin{lemma}[Canonical forms]\label{lem:canonical}
Let $v$ be a closed value.
\begin{enumerate}
\item If $\emptyset\vdash v:\bool$ then $v\in\{\ty{true},\ty{false}\}$.
\item If $\emptyset\vdash v:\num$ then $v=\ul r$ for some $r\in\RR$.
\item If $\emptyset\vdash v:t\arr u$ then $v=\lambda x{:}t.\,e$.
\end{enumerate}
\end{lemma}
\begin{proof}
By case analysis on the grammar of values and inversion of typing.
A value is either $\ul r$, $\ty{true}$, $\ty{false}$, or $\lambda x{:}t.\,e$.
Numerals are typed only by \textsc{T-Num} as $\num$; booleans only by \textsc{T-BoolT/F} as $\bool$; and the only rule concluding an arrow type is \textsc{T-Lam}.
Closedness rules out variables. The three items follow.
\end{proof}

\begin{lemma}[Typed replacement]\label{lem:frames}
All under empty context.
\begin{enumerate}
\item If $\emptyset\vdash e_1\,e_2:u$ and $\emptyset\vdash e_1':t\arr u$ then $\emptyset\vdash e_1'\,e_2:u$.
\item If $\emptyset\vdash v\,e_2:u$ and $\emptyset\vdash e_2':t$ then $\emptyset\vdash v\,e_2':u$.
\item If $\emptyset\vdash \ty{if}\;e_1\;\ty{then}\;e_2\;\ty{else}\;e_3:t$ and $\emptyset\vdash e_1':\bool$ then $\emptyset\vdash \ty{if}\;e_1'\;\ty{then}\;e_2\;\ty{else}\;e_3:t$.
\item If $\emptyset\vdash e_1{<}e_2:\bool$ and $\emptyset\vdash e_1':\num$ then $\emptyset\vdash e_1'{<}e_2:\bool$.
\item If $\emptyset\vdash e_1{<}e_2:\bool$ and $\emptyset\vdash e_2':\num$ then $\emptyset\vdash e_1{<}e_2':\bool$.
\end{enumerate}
\end{lemma}
\begin{proof}
Each item is a one-step inversion on the given typing and reconstruction with the indicated premise.
\end{proof}

\section*{Well-definedness of the semantics}

\begin{proposition}[Bind typing principle]\label{prop:bind-typing}
Let $A,B$ be measurable spaces.
If $\mu\in \D(A)$ and $f:A\to \D(B)$, then $\mu\bind f\in \D(B)$.
In particular, if $f(a)=\dirac(b_a)$ with $b_a\in B$ for all $a\in A$, then $\mu\bind f\in \D(B)$.
\end{proposition}

\begin{theorem}[Totality (progress-like)]\label{thm:total}
For every closed, well-typed $e\in \E{t}$ exactly one defining clause of $\sem{e}$ applies, so $\sem{e}$ is defined.
\end{theorem}
\begin{proof}
By induction on the derivation of the closed typing judgment $\emptyset\vdash e:t$.
We do a case analysis on the last rule used.

T-Num/T-BoolT/T-BoolF/T-Unif/T-Lam: In each case $e$ is a value (numeral, boolean, or abstraction) or a random source. The corresponding Dirac/uniform clause applies, defining $\sem{e}$.

T-App: We have $\emptyset\vdash e_1:t\arr u$ and $\emptyset\vdash e_2:t$. By the IH for the premises, $\sem{e_1}$ and $\sem{e_2}$ are defined.
If $\neg\val(e_1)$ the left-congruence-by-bind clause uses $\sem{e_1}\bind f$, which is defined since $\sem{e_1}$ is defined by IH. Otherwise $\val(e_1)$, and by Lemma~\ref{lem:canonical} $e_1=\lambda x{:}t.\,e$. If $\val(e_2)$ the $\beta$-clause applies; if $\neg\val(e_2)$ the right-congruence-by-bind clause uses $\sem{e_2}\bind f$, which is defined by the IH on the second premise.

T-If: We have $\emptyset\vdash e_1:\bool$, $\emptyset\vdash e_2:t$, $\emptyset\vdash e_3:t$. If $\val(e_1)$, Lemma~\ref{lem:canonical} gives $e_1\in\{\ty{true},\ty{false}\}$ and the corresponding Dirac clause applies. Otherwise $\neg\val(e_1)$ and the guard-congruence-by-bind clause uses $\sem{e_1}\bind f$, which is defined by the IH on the premise typing $\emptyset\vdash e_1:\bool$.

T-Lt: We have $\emptyset\vdash e_1:\num$ and $\emptyset\vdash e_2:\num$. By IH, $\sem{e_1}$ and $\sem{e_2}$ are defined. If both are values, Lemma~\ref{lem:canonical} gives numerals so one of the two Dirac clauses returns a closed boolean. If $\neg\val(e_1)$, the left-congruence-by-bind clause uses $\sem{e_1}\bind f$, defined by IH. Otherwise $\val(e_1)$ and $\neg\val(e_2)$, so the right-congruence-by-bind clause uses $\sem{e_2}\bind f$, defined by IH.

T-Var: Impossible under the empty context.

In all cases the side conditions of the clauses are mutually exclusive, so exactly one clause applies; and in each congruence case the IH guarantees the bind's input is defined; hence $\sem{e}$ is defined.
\end{proof}

\begin{theorem}[Well-formedness (preservation-like)]\label{thm:pres}
For every $e\in \E{t}$, $\sem{e}\in \D(\E{t})$.
\end{theorem}
\begin{proof}[By induction on the typing derivation]
We proceed by induction on the derivation of $\emptyset\vdash e:t$, doing a case split on its last rule.

T-Num/T-BoolT/T-BoolF/T-Lam: $e$ is a value of type $t$, so the corresponding Dirac clause gives $\dirac(e)\in \D(\E{t})$ since $e\in \E{t}$.

T-Unif: The clause is $\sem[\num]{\ty{unif}(a,b)}=\iota_{*}(U[a,b])\in \D(\E{\num})$.

T-App: We have $\emptyset\vdash e_1:t\arr u$ and $\emptyset\vdash e_2:t$. By the IH, $\sem{e_1}\in\D(\E{t\arr u})$ and $\sem{e_2}\in\D(\E{t})$.
- If $\neg\val(e_1)$, the semantics uses $\sem{e_1}\bind f$ with $f(e_1')=\dirac(e_1'\,e_2)$. By Lemma~\ref{lem:frames}(1), $e_1'\in\E{t\arr u}$ implies $e_1'\,e_2\in\E{u}$, so $f:\E{t\arr u}\to \D(\E{u})$. Now apply Proposition~\ref{prop:bind-typing} with $\mu=\sem{e_1}$ to conclude $\sem{e_1}\bind f\in \D(\E{u})$.
- Otherwise $\val(e_1)$, hence by Lemma~\ref{lem:canonical} $e_1=\lambda x{:}t.\,e$. If $\val(e_2)$, then Lemma~\ref{lem:subst} gives $\emptyset\vdash e[v/x]:u$, so the $\beta$-clause $\dirac(e[v/x])$ lies in $\D(\E{u})$. If $\neg\val(e_2)$, use the right bind clause with $f(e_2')=\dirac(e_1\,e_2')$, which is typed by Lemma~\ref{lem:frames}(2); with the IH $\sem{e_2}\in\D(\E{t})$, Proposition~\ref{prop:bind-typing} yields $\sem{e_2}\bind f\in \D(\E{u})$.

T-If: We have $\emptyset\vdash e_1:\bool$, $\emptyset\vdash e_2:t$, $\emptyset\vdash e_3:t$. If $\val(e_1)$ then by Lemma~\ref{lem:canonical} $e_1\in\{\ty{true},\ty{false}\}$ and the corresponding Dirac clause returns an element of $\D(\E{t})$. Otherwise $\neg\val(e_1)$ and the guard bind clause uses $f(e_1')=\dirac(\ty{if}\;e_1'\;\ty{then}\;e_2\;\ty{else}\;e_3)$, which maps $\E{\bool}\to \D(\E{t})$ by Lemma~\ref{lem:frames}(3); with the IH $\sem{e_1}\in\D(\E{\bool})$, conclude by Proposition~\ref{prop:bind-typing}.

T-Lt: We have $\emptyset\vdash e_1:\num$ and $\emptyset\vdash e_2:\num$. If both are values, Lemma~\ref{lem:canonical} gives numerals so one of the two Dirac clauses returns a closed boolean, hence an element of $\E{\bool}$. If $\neg\val(e_1)$, use the left bind clause with $f(e_1')=\dirac(e_1'{<}e_2)$, which maps $\E{\num}\to \D(\E{\bool})$ by Lemma~\ref{lem:frames}(4); with the IH $\sem{e_1}\in\D(\E{\num})$, Proposition~\ref{prop:bind-typing} gives well-formedness. Otherwise $\val(e_1)$ and $\neg\val(e_2)$; use the right bind clause (typed by Lemma~\ref{lem:frames}(5)) together with the IH $\sem{e_2}\in\D(\E{\num})$ and apply Proposition~\ref{prop:bind-typing}.

T-Var: Impossible in the empty context.

Thus in every typing case, $\sem{e}\in \D(\E{t})$.
\end{proof}

\section*{Measurable structure on closed expressions}

The set $\E{t}$ is uncountable due to real numerals. We now construct a standard Borel measurable space on $\E{t}$ by factoring out numerals.

\subsection*{Templates with numeric holes}
\begin{definition}[Templates]
Fix a countable set of variables. Introduce \emph{numeric holes} $\bullet_i$ ($i\in\N$), each intended to be filled by a real numeral.
Templates $\tau$ are generated by
\[
\begin{array}{rcl}
\tau &::=& x \mid \lambda x{:}t.\,\tau \mid \tau\ \tau \mid \ty{if}\;\tau\;\ty{then}\;\tau\;\ty{else}\;\tau \\
&& \mid\ \tau{<}\tau \mid \ty{true}\mid \ty{false}\mid \bullet_i \mid \ty{unif}(\bullet_i,\bullet_j)
\end{array}
\]
with the understanding that $\bullet_i$ has type $\num$.
Let $\holes(\tau)\subseteq\N$ be the finite set of indices used in $\tau$, and let $n(\tau)=|\holes(\tau)|$.
\end{definition}

\begin{lemma}[Countability]\label{lem:countable-templates}
The set of templates is countable. For each fixed $t$, the set $\Tt_t$ of \emph{closed, typable} templates of type $t$ (under the rule that every $\bullet_i$ has type $\num$) is countable.
\end{lemma}
\begin{proof}
Templates are finite trees over a finite operator alphabet with node labels drawn from a countable set (variables and indices). Such trees form a countable set. The subset of closed, typable templates of type $t$ is a subset, hence countable.
\end{proof}

\begin{definition}[Plugging and shaping]\label{def:plug-shape}
Fix a total order on the occurrences of holes in a template (e.g.\ left-to-right). For a template $\tau$ with $n(\tau)=k$, define
\[
\plug(\tau,\vec r)\ \in \E{t}\qquad(\vec r\in \RR^{k})
\]
as the closed term obtained by replacing the $j$-th hole occurrence by the numeral $\ul{r_j}$.

Define the \emph{shape map} $\shape:\E{t}\to \bigsqcup_{\tau\in\Tt_t}\RR^{n(\tau)}$ recursively by scanning $e$ left-to-right and replacing each numeral $\ul r$ and each $\ty{unif}(a,b)$'s numeric arguments $a,b$ by fresh consecutive holes, recording the encountered reals into a vector $\vec r$; the result is a pair $(\tau,\vec r)$ with $\tau\in \Tt_t$.
\end{definition}

\begin{proposition}[Bijection]\label{prop:bij}
For each $t$,
\[
J_t:\ \bigsqcup_{\tau\in\Tt_t}\RR^{n(\tau)} \longrightarrow \E{t},\qquad J_t(\tau,\vec r)=\plug(\tau,\vec r)
\]
is a bijection with inverse $\shape$.
\end{proposition}
\begin{proof}
By construction, $\plug(\shape(e))=e$ for all $e\in \E{t}$.
Conversely, if $\shape(\plug(\tau,\vec r))=(\tau',\vec r')$, the canonical left-to-right numbering of holes yields $\tau'=\tau$ and $\vec r'=\vec r$.
Disjointness across different $\tau$ is ensured by the disjoint union.
\end{proof}

\subsection*{Measurable space}
Equip each $\RR^{k}$ with its Borel $\sigma$-algebra $\B(\RR^{k})$ and the disjoint union $\bigsqcup_{\tau\in\Tt_t}\RR^{n(\tau)}$ with the sum $\sigma$-algebra
\[
\Sigma^\sqcup_t \;=\; \bigl\{\, \bigcup_{\tau\in\Tt_t} \{\tau\}\times A_\tau \ \big|\ A_\tau\in \B(\RR^{n(\tau)}) \,\bigr\}.
\]
\begin{definition}[Measurable space of closed terms]\label{def:measspace}
Define $\Sigma_t$ on $\E{t}$ as the pushforward $\sigma$-algebra $\,\Sigma_t \triangleq (J_t)_{*}(\Sigma^\sqcup_t)$.
Then $J_t$ is a measurable isomorphism between $(\bigsqcup_{\tau}\RR^{n(\tau)},\Sigma^\sqcup_t)$ and $(\E{t},\Sigma_t)$.
\end{definition}

\begin{proposition}[Standard Borel]\label{prop:standard-borel}
$(\E{t},\Sigma_t)$ is a standard Borel space. Consequently, subprobability measures on $\E{t}$ form the carrier of the Giry monad; $\dirac$ and $\bind$ are well-defined, and the maps used in the semantics are measurable.
\end{proposition}

\section*{Corollary (one-step Markov kernel)}
By Theorems~\ref{thm:total} and~\ref{thm:pres}, and the measurable-space construction above, for each $t$ the mapping $\sem[t]{\cdot}:(\E{t},\Sigma_t)\to \D(\E{t},\Sigma_t)$ is a total, type-preserving one-step Markov kernel.

\end{document}
